\documentclass[11pt,a4paper]{article}
\usepackage{templates/modeling-notes}
\hypersetup{pdftitle={Chapter 1 -- Modeling optimization problems}}
\title{\textbf{Mathematical Modeling}\\[2mm]
\large Chapter 1: Modeling Optimization Problems}
\author{Personal study notes}
\date{}
\begin{document}
\maketitle
\tableofcontents

\section{What mathematical programming means}\label{sec:meaning}
Here, \emph{mathematical programming} means mathematical optimization: selecting
the best feasible decision according to a specified criterion. The word
``programming'' refers to planning; an optimization problem is a mathematical
object even before any computer program is written.

A model has three central ingredients:
\begin{enumerate}[leftmargin=1.8em]
  \item \textbf{Decision variables}: quantities that may be chosen.
  \item \textbf{Constraints}: conditions that an acceptable decision must satisfy.
  \item \textbf{Objective function}: a numerical criterion to minimize or maximize.
\end{enumerate}
Parameters, such as available resources or prices, are fixed data for one model
instance. A parameter can change between instances without becoming a decision
variable.

\section{A concrete model: five products and three resources}\label{sec:factory}
Suppose a factory chooses weekly amounts $x_1,\ldots,x_5$ of five products.
Treat these amounts as nonnegative real quantities: this is a continuous model
whose divisibility assumption must be justified for the actual products.
Machine use and profit per unit are fixed and additive:
\begin{center}
  \begin{tabular}{lrrrrr}\toprule
    Product & 1 & 2 & 3 & 4 & 5\\\midrule
    Profit (PLN/unit) & 550 & 600 & 350 & 400 & 200\\
    Grinding (h/unit) & 12 & 20 & 0 & 25 & 15\\
    Drilling (h/unit) & 10 & 8 & 16 & 0 & 0\\
    Technician work (h/unit) & 20 & 20 & 20 & 20 & 20\\\bottomrule
  \end{tabular}
\end{center}
Each machine and technician works $6\cdot8=48$ hours per week. Two drilling
machines supply $96$ machine-hours, and eight technicians supply $384$
person-hours. Let $G$ denote available grinding hours; three machines give
$G=144$, whereas four give $G=192$. Assume capacity can be pooled, with no
additional scheduling, setup, or demand restrictions. Each coefficient times
$x_j$ is a required number of hours, so the model is
\begin{equation}\label{eq:factory}
  \begin{aligned}
    \max\quad &550x_1+600x_2+350x_3+400x_4+200x_5\\
    \text{subject to}\quad
    &12x_1+20x_2+25x_4+15x_5\leq G,\\
    &10x_1+8x_2+16x_3\leq96,\\
    &20(x_1+x_2+x_3+x_4+x_5)\leq384,\qquad x\geq0.
  \end{aligned}
\end{equation}
Resource limits give upper bounds, while $x_j\geq0$ prevents negative production.
The objective counts total profit, using the stated unit profits. Choosing a
single product by its profit per unit alone ignores the different resources it uses.
For a general indexed model, let row $i$ represent a resource and column $j$
a product: $a_{ij}$ has units of resource $i$ per unit of product $j$.
Then row $i$ is $\sum_j a_{ij}x_j\leq b_i$, or $Ax\leq b$.
Rows combine consumption of one resource across all products; columns describe
all resources used by one product.

\subsection{A hand solution with a reason it is optimal}
For $G=144$ or $192$, consider
\[
  x_1=\frac{48}{5},\quad x_4=\frac{G-576/5}{25},\quad x_2=x_3=x_5=0.
\]
Grinding and drilling are used exactly; technician work is below $384$ hours.
To compare this plan with \emph{every} feasible plan, multiply the grinding
inequality by $16$ and the drilling inequality by $179/5$. Their combined
left side has coefficients
\[
  (550,\ 3032/5,\ 2864/5,\ 400,\ 240),
\]
each at least the corresponding profit coefficient. Since $x\geq0$,
\[
  550x_1+600x_2+350x_3+400x_4+200x_5
  \leq16G+\frac{179}{5}\,96.
\]
The proposed plan attains this upper bound, so it is optimal. This weighted
resource argument is an \emph{optimality certificate}: the bound and its
attainment establish the answer without trusting a solver.
\begin{center}
  \begin{tabular}{rrrr}\toprule
    $G$ (h) & $x_1$ & $x_4$ & Optimal profit (PLN)\\\midrule
    $144$ & $48/5$ & $144/125$ & $28704/5$\\
    $192$ & $48/5$ & $384/125$ & $32544/5$\\\bottomrule
  \end{tabular}
\end{center}
The strictly larger bound coefficients for products 2, 3, and 5 force their
amounts to zero at an optimum; equality in the two resource bounds then fixes
$x_1$ and $x_4$. Each continuous optimum is unique. Requiring whole products
would change the model; rounding these quantities does not establish an integer optimum.

\section{The general constrained problem}\label{sec:general}
Let $x=(x_1,\ldots,x_n)^T\in\R^n$ be the decision vector. A constrained
minimization problem is
\begin{equation}\label{eq:general}
  \min_{x\in\Omega} f(x),\qquad
  \Omega=\{x\in\R^n:g_i(x)\leq0\ (i\in I),\ h_j(x)=0\ (j\in E)\}.
\end{equation}
Here $f$ is the objective, $g_i$ are inequality-constraint functions, and $h_j$
are equality-constraint functions; $I$ and $E$ are finite index sets. Bounds and
sign restrictions must also be included in $\Omega$ when required.

\begin{definition}[Feasibility and global optimality]\label{def:optimal}
  A point $x$ is \emph{feasible} if $x\in\Omega$. A feasible point $x^*$ is a
  \emph{global minimizer} if $f(x^*)\leq f(x)$ for every $x\in\Omega$.
  The value $f(x^*)$ is the optimal objective value.
\end{definition}
An empty feasible set means that the requirements cannot be met simultaneously.
Feasibility alone does not imply optimality. The optimal decision $x^*$ and its
objective value $f(x^*)$ are different objects, and multiple decisions may have
the same value.

\begin{remark}[Minimum versus infimum]
  A finite lower bound need not be attained in a general optimization problem.
  For example, $\inf_{x>0}x=0$, but no feasible point has value zero.
  A solver's output therefore requires both feasibility and an appropriate
  optimality or attainment argument.
\end{remark}

\section{From a real question to a mathematical model}\label{sec:workflow}
\begin{enumerate}[leftmargin=1.8em]
  \item \textbf{Choose variables and construct the model.} What can actually be
        controlled? What is given? State each variable's meaning, domain, and unit.
  \item \textbf{Specify constraints.} Translate capacities, balances, required
        levels, and physical limits into equations or inequalities. Check their direction.
  \item \textbf{Specify the objective.} State what ``best'' means, whether the
        criterion is minimized or maximized, and which units measure it.
  \item \textbf{Select an algorithm.} Match the method to the model's structure:
        linearity, smoothness, convexity, dimensions, and continuous or discrete variables.
  \item \textbf{Find and verify a solution.} Check feasibility and theoretical
        optimality conditions, then interpret the answer in the original setting.
  \item \textbf{Perform post-optimal analysis.} Investigate how the answer depends
        on data and assumptions.
\end{enumerate} A correct solution to an inaccurate
model can still be a poor real-world decision.

\subsection{Units and assumptions}
In a resource constraint, a coefficient measured in kg/item times a variable
measured in items produces kg; its right-hand side must also be kg. Price in
PLN/item times quantity in items produces revenue in PLN.

State whether quantities may be fractional. A count of indivisible objects
requires integrality; permitting real values gives a continuous relaxation.
State assumptions such as fixed coefficients, additive contributions, and
unchanging prices. These are claims about the setting, not facts guaranteed by
an optimization algorithm.

\section{How optimization problems are classified}\label{sec:classification}
\begin{description}[leftmargin=2em,style=nextline]
  \item[Constraints] Unconstrained or constrained optimization.
  \item[Variables over time] Static decisions or decisions evolving over time.
        Dynamic problems include control and functional optimization;
        this is broader than the algorithmic technique called dynamic programming.
  \item[Number of objectives] Single-objective or multiobjective optimization.
        Several competing criteria require an explicit way to compare trade-offs.
  \item[Function structure] Linear, nonlinear, quadratic, or convex optimization.
        Smoothness is another distinction; a nonlinear function may be smooth or nonsmooth.
  \item[Allowed values] Continuous, discrete, integer, or binary variables.
\end{description}
A linear program has a linear objective, linear
constraints, and continuous variables. Adding integrality produces an integer
linear program. Quadratic programming does not automatically mean convexity.
For $f(x)=\tfrac12x^TQx+c^Tx$ with symmetric $Q$, convexity requires $Q$ to be
positive semidefinite; a convex optimization problem also needs a convex feasible set.

\section{Optimization roadmap}\label{sec:roadmap}
\begin{enumerate}[leftmargin=1.8em]
  \item \textbf{Linear programming}: modeling examples, standard and canonical
        forms, graphical solutions, bases, simplex methods, duality, post-optimal analysis,
        and interior-point methods.
  \item \textbf{Unconstrained nonlinear optimization}: local optimality conditions,
        iterative methods, convergence, line searches, steepest descent, Newton and
        quasi-Newton methods, conjugate gradients, and direct search.
  \item \textbf{Constrained nonlinear optimization}: Lagrange multipliers,
        Karush--Kuhn--Tucker conditions, quadratic and convex problems, penalty and barrier
        methods, and interior-point methods.
\end{enumerate}
The simplex method for linear programming and the Nelder--Mead simplex method
are different algorithms. Their shared name does not make them interchangeable.

% Keep the closing study sections together; avoid a crowded page 3 and sparse page 4.
\newpage
\section{MATLAB refresher scope}\label{sec:matlab}
Core MATLAB skills include vectors and matrices,
programming constructs, built-in functions (including numerical methods), scripts,
function files, parameter passing (including passing functions), and basic plots.
Optimization Toolbox provides solvers for comparison with implemented algorithms.

\section{Questions for revision}\label{sec:revision}
\begin{enumerate}[leftmargin=1.8em]
  \item Which quantities in a model are decisions and which are parameters?
  \item What distinguishes a feasible point, an optimal point, and an optimal value?
  \item How would changing an inequality's direction change the model?
  \item Why can fractional counts be mathematically feasible but physically unusable?
  \item What must be checked after an algorithm returns a candidate solution?
\end{enumerate}

\section{Sources and editorial remarks}\label{sec:sources}
\noindent\textbf{Primary source:} \texttt{s4\_w1.pdf}, PDF pages 1--6.\\
\textbf{Source context:} Folder labelled 2020; the slides contain a 2019/2020
schedule. Exact lecture date is not established.\\
\textbf{Current course alignment:} The supplied 2026--27 lecture
\texttt{ProgMat\_w1.pdf}, pp.~2--8, supports the modeling example in
Section~\ref{sec:factory}; the remaining topics are mapped in the
\href{../../docs/course-map-2026-27.md}{current-course map}.
The historical roadmap and MATLAB scope
are not confirmed current syllabus or assessment instructions.\\
\textbf{Editorial scope:} English synthesis of the mathematical content.
Definitions and explanations expand the brief slides; historical schedules,
assessment rules, and contact details are not current course instructions.

\paragraph{Current production data and verification.}
\texttt{ProgMat\_w1.pdf}, pp.~2--7, supplies the five-product data.
Page~4 states four grinding machines, while pp.~6--7 use three and $144$ hours.
Section~\ref{sec:factory} retains capacity as $G$ and explicitly solves both
readings. The weighted bound and uniqueness argument are added explanations.
Page~8 indexes consumption by product first and resource second. The resource-row
matrix convention in Section~\ref{sec:factory} is its transpose; changing index
names consistently preserves the resource sums.
The actual companion is
\href{run:matlab/exercises/modeling_patterns_checks.m}{modeling-patterns MATLAB script} at \path{matlab/exercises/modeling_patterns_checks.m};
it checks the same rational data with absolute tolerance $10^{-10}$.
The new results were checked independently in exact rational arithmetic.
MATLAB execution of this new script is unverified; execution evidence is maintained
in \path{docs/matlab-verification.md}.

Polish equivalents are maintained in \texttt{docs/glossary.md}, including
\emph{zmienna decyzyjna}, \emph{funkcja celu}, \emph{ograniczenie},
\emph{rozwi\k{a}zanie dopuszczalne}, and \emph{rozwi\k{a}zanie optymalne}.
\begin{itemize}[leftmargin=1.6em]
  \item \texttt{s4\_w1.pdf}, pp.~1--2: historical administration and bibliography;
        not reproduced as current instructions.
  \item pp.~3--4: Sections~\ref{sec:meaning}, \ref{sec:general}, \ref{sec:workflow},
        and \ref{sec:classification}, and the start of
        Section~\ref{sec:roadmap}.
  \item pp.~5--6: remaining roadmap in Section~\ref{sec:roadmap} and MATLAB scope in
        Section~\ref{sec:matlab}.
\end{itemize}
Page numbers refer to PDF pages, each containing up to four slides.
\end{document}
